%=============================================================
% Chapitre 0 — Introduction générale
%=============================================================
\chapter*{Introduction générale}
\addcontentsline{toc}{chapter}{Introduction générale}
\markboth{Introduction générale}{}

\section*{Contexte}

Dans un contexte de forte polarisation politique et de montée des incertitudes électorales,
la capacité à anticiper le comportement des électeurs à partir de données socio-économiques
représente un enjeu stratégique majeur pour les acteurs politiques, les instituts de sondage
et les cabinets de conseil en stratégie.

La startup fictive \textbf{Electio-Analytics} nous mandate pour concevoir un système de
prédiction électorale appliqué à la région \textbf{Nouvelle-Aquitaine}, à partir des
indicateurs des recensements INSEE 2016 et 2022. Cette région, avec ses 12 départements,
ses dynamiques économiques contrastées et sa diversité sociologique, constitue un terrain
d'étude riche et représentatif des fractures territoriales françaises.

\section*{Problématique}

\begin{center}
\textit{« Dans quelle mesure les évolutions socio-économiques locales entre 2016 et 2022
permettent-elles de prédire l'orientation politique des électeurs en Nouvelle-Aquitaine,
et comment modéliser cette relation par des algorithmes de machine learning ? »}
\end{center}

\section*{Objectifs du projet}

\begin{enumerate}
  \item \textbf{Collecter et structurer} les données socio-économiques INSEE (emploi,
        logement, population, chômage) pour les 627 cantons de Nouvelle-Aquitaine.
  \item \textbf{Ingénierie des variables} : calculer les variations inter-recensements
        ($\Delta$ features) pour capturer les dynamiques de changement.
  \item \textbf{Construire un pipeline ETL} en 8 phases, avec \textbf{MongoDB} comme socle de stockage central : collecte brute, base préparée, résultats ML et prédictions.
  \item \textbf{Entraîner et évaluer} cinq modèles de machine learning sur une variable
        cible à 5 classes synthétiques.
  \item \textbf{Prédire le comportement électoral} au niveau départemental et régional
        pour la présidentielle 2022.
  \item \textbf{Restituer les résultats} via une interface React 19 interactive.
\end{enumerate}

\section*{Plan du mémoire}

Le présent mémoire s'organise en sept chapitres :

\begin{description}
  \item[Chapitre 1 — Cadre géographique :] présentation de la Nouvelle-Aquitaine,
    ses 12 départements et le contexte électoral 2022.
  \item[Chapitre 2 — Données et indicateurs :] sources, 171 variables,
    construction des $\Delta$-features.
  \item[Chapitre 3 — Architecture et pipeline :] les 8 phases du pipeline ETL,
    MongoDB comme stockage central (collecte brute $\rightarrow$ \texttt{raw\_data},
    base finale $\rightarrow$ \texttt{data\_final}, résultats ML $\rightarrow$ \texttt{predictions}).
  \item[Chapitre 4 — Analyse exploratoire :] EDA, corrélations, distributions.
  \item[Chapitre 5 — Modélisation ML :] cinq algorithmes, validation croisée,
    sélection du meilleur modèle.
  \item[Chapitre 6 — Résultats et interprétation :] prédictions départementales,
    précision à 91,7\%, analyse des erreurs.
  \item[Chapitre 7 — Interface applicative :] dashboard React 19, visualisations
    interactives, architecture technique.
\end{description}
